\begin{textsample}{NEG Dim 2 – 2011\_c5 – Score -3.00 – 2011\_c5/enade\_2011\_eng\_g\_6\_7d2177b9ec.txt}  \label{ex:f2_neg_085}
CAPÍTULO 5 DISTRIBUIÇÃO DOS CONCEITOS 5.1 PANORAMA NACIONAL DA DISTRIBUIÇÃO DOS CONCEITOS A Tabela 5.1 apresenta a quantidade e distribuição de cursos de Engenharia – Grupo VI participantes do ENADE/2011, por faixa de conceito e Grande Região.
A diferença entre os cursos tabulados neste capítulo e no capítulo 2 são os cursos sem conceito, em princípio, aqueles sem alunos concluintes que participassem da prova.
Observando -se os dados da Tabela 5.1, nota -se que, dos 215 cursos participantes, 87 ( 40,4 \% ) classificaram -se com conceito 2, o valor modal.
Este foi também o conceito modal nas regiões Nordeste ( 38,4 \% ), Sudeste ( 46,8 \% ) e Centro-Oeste ( 37,5 \% ).
Nas regiões Norte e Sul o conceito modal foi 3 com, respectivamente, 50,0 \% e 46,4 \% dos cursos da região.
O conceito 3 foi o segundo mais frequente em nível nacional ( 27,4 \%, correspondendo a 59 cursos ) e o conceito 4, o terceiro ( 14,9 \%, correspondendo a 32 cursos ).
Houve, ainda, nove cursos ( 4,2 \% ) que receberam conceito 5 e outros 24 cursos ( 11,2 \% ) que receberam conceito 1.
Dos 215 cursos de Engenharia – Grupo VI quatro ficaram sem conceito ( SC ).
A região Norte participou com 12 cursos ou 5,6 \% do total nacional.
Nesta região, o conceito 3 foi atribuído a seis cursos dos 12 participantes, 50,0 \% do total regional.
Nenhum curso ficou sem conceito.
Os demais foram avaliados com conceito 1 ( dois cursos, correspondendo a 16,7 \% ), conceito 2 ( três cursos correspondendo a 25,0 \% ) e conceito 4 ( um curso correspondendo a 8,3 \% ).
Nenhum curso alcançou o conceito 5.
A região Nordeste participou com 26 cursos ou 12,1 \% do total nacional.
Como já comentado, destes, dez cursos, 38,4 \% em termos regionais, obtiveram conceito 2, o conceito modal para a região.
Os conceitos 3 e 4 foram atribuídos a, respectivamente, oito e quatro cursos ( 30,8 \% e 15,4 \% ).
O conceito 1 foi atribuído a outros quatro cursos ( 15,4 \% ).
Nessa região nenhum dos cursos ficou sem conceito tampouco recebeu o conceito 5.
Dos 126 cursos participantes da região Sudeste ( 58,6 \% do total nacional ), 59 ( 46,8 \% da região ) obtiveram conceito 2, o conceito modal.
O conceito 1 foi atribuído a 14 cursos ( 11,1 \% ) e o conceito 3, a 24 ( 19,0 \% ).
Vinte cursos ( 15,9 \% ) receberam o conceito 4 e seis ( 4,8 \% ) receberam o conceito 5.
Três cursos ficaram sem conceito.
A região Sul também contou com cursos \textbf{distribuídos} em todas as faixas de conceitos.
A predominância do conceito 3 foi de 46,4 \%, correspondentes a 20 dos 43 cursos participantes na região Sul.
O conceito 2 foi atribuído a 12 cursos ( 27,9 \% ) e o conceito 4, a seis cursos ( 14,0 \% ).
O conceito 1 foi recebido por dois cursos ( 4,7 \% ) e o conceito 5, por três cursos ( 7,0 \% ).
Nenhum dos cursos da região Sul ficou sem conceito.
Dos oito cursos participantes na região Centro-Oeste 37,5 \%, correspondendo a três cursos, recebeu conceito 2, o conceito modal.
Os demais cursos foram avaliados com conceitos : 1 ( dois cursos, 25,0 \% ), 3 ( um curso, 12,5 \% ) e 4 ( também um curso, 12,5 \% ).
Nesta região um curso ficou sem conceito e nenhum alcançou o conceito 5. 5.2 CONCEITOS POR CATEGORIA ADMINISTRATIVA E POR GRANDE REGIÃO A Tabela 5.2 apresenta a distribuição dos cursos de Engenharia – Grupo VI participantes do ENADE/2011, por Categoria Administrativa, de acordo com os conceitos por eles alcançados, segundo as Grandes Regiões brasileiras.
Dos 215 cursos participantes, 58 ( 27,0 \% ) eram ministrados em instituições públicas e 157 ( 73,0 \% ), em privadas.
De acordo com as informações da Tabela 5.2, em termos nacionais quatro cursos ficaram sem conceito dois eram públicos e outros dois eram privados.
Dos 58 cursos participantes de IES públicas, o conceito 4 foi o valor modal, atribuído a 23 cursos.
Entre os demais cursos participantes, um curso obteve conceito 1 ( 1,7 \% da categoria ), oito receberam conceito 2 ( 13,8 \% ) e 16 foram avaliados com conceito 3 ( 27,6 \% ).
Os oito cursos ( 13,8 \% ) restantes receberam o conceito 5.
Nesta categoria, como já comentado, dois dos cursos ficaram sem conceito.
Na rede privada, o conceito modal foi 2, com 79 cursos dos 157 da categoria ( 50,3 \% ).
Entre os demais cursos participantes, 23 receberam conceito 1 e 43 o conceito 3.
O conceito 4 foi atribuído a nove cursos e o conceito 5 a apenas um curso.
Nesta categoria administrativa, dois dos cursos ficaram sem conceito.
Na análise por região, observa -se que, na região Norte, as instituições privadas participaram com seis cursos ( 50,0 \% do total regional ), dos quais nenhum ficou sem conceito.
O conceito modal para as instituições privadas na região foi o conceito 3, com três cursos.
Dois cursos receberam o conceito 1 ( 33,3 \% ) e um outro curso recebeu conceito 2 ( 16,7 \% ).
Nenhum curso recebeu os dois conceitos mais altos, 4 e 5.
As instituições públicas participaram com seis cursos ( outros 50,0 \% do total regional ).
O conceito modal para esta Categoria Administrativa, na região Norte, também foi o conceito 3, com três cursos.
Dois cursos receberam o conceito 2 ( 33,3 \% ) e o outro curso ( 16,7 \% ) recebeu o conceito 4.
Nenhum curso ficou sem conceito ou recebeu os conceitos 1 e 5 nesta categoria na região.
Na região Nordeste, a rede privada concentrou 15 dos 26 cursos participantes, equivalentes a 57,7 \% do total da região.
Nesta Categoria Administrativa na região, \textbf{predominaram} os cursos com conceito 2, oito cursos correspondendo a 53,3 \%.
Os demais foram avaliados com conceitos 1 ( quatro cursos correspondendo a 26,7 \% ) e 3 ( três cursos correspondendo a 20,0 \% ).
Nenhum curso recebeu os conceitos 4 ou 5, nem tampouco houve curso sem conceito.
As instituições públicas dessa região participaram com 11 cursos ( 42,3 \% ), dos quais cinco obtiveram o conceito 3, modal.
Os demais foram avaliados com os conceitos 2 ( dois cursos ) e 4 ( quatro cursos ).
Nenhum curso ficou sem conceito, nesta combinação de região e categoria, ou recebeu conceitos 1 ou 5.
Na região Sudeste, a proporção de cursos da rede privada, 81,0 \%, foi a mais \textbf{elevada} dentre as regiões brasileiras, correspondendo a 102 dos 126 cursos participantes.
Nesta categoria, na região Sudeste, o conceito modal foi 2 ( 57 cursos ).
Os demais foram avaliados com conceito 1 ( 14 cursos ), conceito 3 ( 20 cursos ), conceito 4 ( oito cursos ), e conceito 5 ( um curso ).
Nesta combinação de Categoria Administrativa e Grande Região, dois cursos não receberam conceito.
Entre os 24 cursos oferecidos em instituições públicas na região Sudeste, a categoria modal foi o conceito 4 ( 12 cursos ).
Os demais foram avaliados com os conceitos : 2 ( dois cursos ), 3 ( quatro cursos ) e 5 ( cinco cursos ).
Na região Sudeste nesta categoria, um curso ficou sem conceito.
As instituições privadas concentraram 29 dos 43 cursos participantes da região Sul, 67,4 \% do total regional.
Desses, 16 ficaram com conceito 3, o conceito modal.
Os demais foram avaliados com conceito 1 ( dois cursos ), conceito 2 ( dez cursos ) e conceito 4 ( um curso ).
Nesta combinação de Categoria Administrativa e Grande Região, nenhum curso ficou sem receber conceito ou alcançou o conceito 5.
As instituições públicas na região Sul participaram com 14 cursos ( 32,6 \% ), dos quais nenhum ficou sem conceito.
O conceito modal foi o 4, com cinco cursos.
Os demais cursos receberam conceitos : 2 ( dois cursos ), 3 ( quatro cursos ) e 5 ( três cursos ).
Nenhum curso foi avaliado com conceito 1.
Na região Centro-Oeste, cinco dos oito cursos participantes eram de instituições privadas ( 62,5 \% em termos regionais ).
Destes, três obtiveram conceito 2, conceito modal.
Os demais receberam conceito 1 e 3 ( um curso cada ).
Nenhum curso ficou sem conceito, nem tampouco recebeu os conceitos 4 e 5.
Dos dois cursos de instituições públicas que receberam conceito, um obteve conceito 1 e outro obteve conceito 4.
Nesta região, um curso de IES pública ficou sem conceito. 5.3 CONCEITOS POR ORGANIZAÇÃO ACADÊMICA E POR GRANDE REGIÃO Na Tabela 5.3 encontra -se a distribuição dos conceitos atribuídos aos cursos participantes do ENADE/2011 na Área de Engenharia – Grupo VI, por Organização Acadêmica, segundo as Grandes Regiões brasileiras.
Dos 215 cursos participantes, 121 eram oferecidos em Universidades, 40 em Centros Universitários e os demais 54 em Faculdades.
Esta distribuição corresponde a, respectivamente, 56,3 \%, 18,6 \% e 25,1 \% dos cursos.
De acordo com os dados apresentados, oito dos nove cursos avaliados com conceito 5 eram vinculados a Universidades.
Dois dos cursos oferecidos em Universidades ficaram sem conceito.
Esse tipo de Organização Acadêmica teve o conceito 3 como modal, com 41 cursos.
Os demais cursos avaliados receberam os conceitos : 1 ( oito cursos ), 2 ( 34 cursos ), 4 ( 28 cursos ) e 5 ( oito cursos, como já mencionado ).
Entre os 40 cursos oferecidos em Centros Universitários, o conceito modal foi 2, com 19 cursos.
Neste tipo de organização acadêmica um curso ficou sem conceito e nenhum recebeu o conceito 5.
Os outros cursos neste tipo de Organização Acadêmica receberam os conceitos : 1 ( sete cursos ), 3 ( 11 cursos ) e 4 ( dois cursos ).
Nas Faculdades, um dos 54 cursos ficou sem conceito e 34 receberam o conceito modal 2.
Dos demais cursos neste tipo de Organização Acadêmica, nove receberam conceito 1, sete receberam o conceito 3, dois o conceito 4 e um curso o conceito 5.
Nenhum curso ficou sem conceito e receberam todos os conceitos.
Considerando -se separadamente as regiões brasileiras, verifica -se que na região Norte as Universidades concentraram oito dos 12 cursos participantes.
Dos cursos oferecidos em Universidades, nenhum ficou sem conceito ou recebeu conceitos 1 ou 5.
Dois cursos receberam o conceito 2, cinco cursos receberam o conceito modal, 3, e um curso recebeu o conceito 4.
Os Centros Universitários da região Norte foram representados por dois cursos, nenhum sem conceito.
Um recebeu o conceito 1 e um outro o conceito 2.
As Faculdades também participaram com dois cursos na região Norte.
Os dois cursos foram avaliados com conceito 1 e 3 ( um curso cada ).
Na região Nordeste, as Universidades participaram com 15 dos 26 cursos na Área de Engenharia – Grupo VI da região.
Nenhum dos cursos oferecidos em Universidades no Nordeste ficou sem conceito.
O conceito modal foi 3, com seis cursos.
Os demais receberam conceito 1 ( um curso ) e os conceitos 2 e 4 receberam quatro cursos cada.
Os Centros Universitários contaram com três cursos participantes na região Nordeste, sendo dois com conceito 2 e um curso com conceito 1.
As Faculdades foram representadas por oito cursos na região Nordeste e todos receberam conceito.
O conceito modal foi 2, com quatro cursos.
Os quatro cursos restantes se dividiram igualmente entre os conceitos 1 e 3.
Na região Sudeste, as Universidades concentraram 66 dos 126 cursos da região ( 52,4 \% ).
Entre os cursos oferecidos em Universidades na região, o conceito modal foi 2, com 21 cursos e um curso ficou sem conceito.
Os demais cursos receberam os conceitos 1 ( seis cursos ), 3 ( 17 cursos ), 4 ( 16 cursos ) e 5 ( cinco cursos ).
Os Centros Universitários participaram com 27 cursos na região Sudeste, dos quais 13 obtiveram conceito modal, 2, e um ficou sem conceito.
Os demais receberam os conceitos : 1 ( quatro cursos ), 3 ( sete cursos ) e 4 ( dois cursos ).
Nenhum curso nesta combinação de organização acadêmica e região alcançou o conceito 5.
As Faculdades foram representadas por 33 cursos na região Sudeste, que se \textbf{distribuíram} nos conceitos 1 ( quatro cursos ), 2 ( 25 cursos, conceito modal ), 4 ( dois cursos ) e 5 ( um curso ).
Um curso ficou sem conceito e nenhum recebeu o conceito 3.
Dos 43 cursos da região Sul, 26 eram de Universidades, para os quais o conceito modal foi 3, com 12 cursos.
Nesse tipo de organização, nenhum dos cursos ficou sem conceito e os demais receberam os conceitos : 2 ( cinco cursos ), 4 ( seis cursos ) e 5 ( três cursos ).
Os Centros Universitários da região Sul tiveram quatro dos oito cursos participantes no conceito modal, 3.
Os outros cursos receberam conceitos 1 ( um curso ) e 2 ( três cursos ).
Foram nove os cursos vinculados a Faculdades na região Sul, nenhum destes ficou sem conceito.
Um curso recebeu o conceito 1, quatro cursos obtiveram conceito 2 e outros quatro o conceito 3, conceitos modais.
Na região Centro-Oeste, seis dos oito cursos eram de Universidades.
Nesse tipo de organização, um curso ficou sem conceito e o conceito modal foi 2, com dois cursos.
Os demais obtiveram os conceitos 1, 3 e 4 ( um curso cada ).
Os Centros Universitários na região Centro-Oeste não participaram do ENADE/2011.
Dos dois cursos oferecidos em Faculdades na região Centro-Oeste, nenhum ficou sem conceito, um curso recebeu conceito 1 e o outro curso conceito 2.

% matched lemmas: distribuir, elevado, predominar
\end{textsample}
